\bookchapter{ภาคผนวก A: ขั้นตอนแกะกล่องมือถือใหม่ฉบับ พี่เลี้ยงดูทุกขั้น}

\booksection{แกะกล่องตรวจเครื่อง สามขั้นตอน:}

\begin{itemize}
  \item \textbf{ดูภายนอก}: ตรวจตัวเครื่อง หน้าจอว่ามีรอยขีดข่วน รอยกระแทกไหม
  \item \textbf{เทียบข้อมูล}: ตรวจว่า IMEI/หมายเลขซีเรียลบนกล่อง ตรงกับบนใบเสร็จไหม
  \item \textbf{ลองฟังก์ชัน}: โทรออกทดสอบหูฟังไมค์ ถ่ายรูปสักใบตรวจกล้อง
\end{itemize}

\booksection{การเปิดใช้งานและตั้งค่าเริ่มต้น จุดที่ต้องเน้น:}

\begin{itemize}
  \item \textbf{เครือข่าย}: ต่อ Wi-Fi ที่นิ่งให้เรียบร้อยก่อน
  \item \textbf{บัญชี}: ล็อกอิน Apple ID, บัญชี Huawei, บัญชี Xiaomi ฯลฯ ของคุณ
  \item \textbf{การตั้งค่าพื้นฐาน}: เปิด "ค้นหาโทรศัพท์ของฉัน" ตั้งรหัสล็อกหน้าจอที่น่าเชื่อถือหนึ่งอัน
  \item \textbf{ตั้งค่า eSIM}: ถ้ามือถือคุณรองรับ eSIM (เช่น iPhone หรือฟลากชิปแอนดรอยด์บางรุ่น) เพิ่มแพ็กเกจ eSIM ได้ในการตั้งค่าโดยไม่ต้องเสียบซิมจริง ไปต่างประเทศ eSIM สะดวกพิเศษ ซื้อแพ็กเกจ eSIM ของผู้ให้บริการท้องถิ่นได้ทันที
  \item \textbf{การย้ายข้อมูล}: ตอนย้ายข้อมูลจากเครื่องเก่า ให้แน่ใจว่าสองเครื่องมีแบตเพียงพอ iPhone ใช้ "เริ่มต้นด่วน" แอนดรอยด์ใช้แอปย้ายเครื่องทางการ
\end{itemize}

\booksection{ติดตั้งแอป ความปลอดภัยมาก่อน:}

\begin{itemize}
  \item \textbf{เลือกสโตร์ทางการก่อน}: แอปเปิลใช้ App Store แอนดรอยด์ใช้สโตร์ในตัวของเครื่อง
  \item \textbf{อย่าให้สิทธิ์ฟุ่มเฟือย}: คำขอสิทธิ์ที่เด้งขึ้นตอนติดตั้ง คิดให้ดีว่า "แอปนี้จำเป็นจริงหรือ?"
\end{itemize}
